% =============================================================================
% Materiał i metody. Plik edytowany ręcznie -- skrypt analiza_odpowiedzi_naukowa.py go NIE nadpisuje.
% Liczby: definicje.tex; tabele: \TabDir (CSV); rysunki: \FigDir; literatura: references.bib.
% =============================================================================
\section{Materiał i metody}

\subsection{Środowisko i program szkoleniowy}
Program prowadzono w banku. Jego treść stanowiły pytania ogólne pogrupowane w cztery kategorie:
wymogi regulacyjne, produkty bankowe, umiejętności miękkie oraz moduł dla kadry kierowniczej.
Każde pytanie było pytaniem wielokrotnego wyboru z trzema do pięciu opcjami odpowiedzi, z których
dokładnie jedna była poprawna, i zawierało wyjaśnienie poprawnej odpowiedzi. Aplikacja szkoleniowa
uruchamiała się automatycznie na ekranach komputerów pracowników i wyświetlała od jednego do trzech
pytań dziennie w godzinach 10:00--14:00; po każdej odpowiedzi pokazywano poprawną odpowiedź wraz
z wyjaśnieniem. To samo pytanie mogło zostać ponownie zadane temu samemu pracownikowi nie wcześniej
niż siedem dni po poprzedniej odpowiedzi, a harmonogram powtórek nie zależał od tego, czy odpowiedź
była poprawna. Na czas przeglądu i aktualizacji pytań program był zawieszany.
O udziale w programie decydował Departament Digitalizacji i Procesów banku, który był również autorem
pytań. Pytania prezentowano w modułach tematycznych (mediana \OptionsMedian{} zaobserwowanych opcji
odpowiedzi na pytanie). W analizowanych danych pracownicy odpowiadali na medianę \DailyMedian{}
pytań w dniu aktywności, \WeekdayPct\% odpowiedzi udzielono w dni robocze, a to samo pytanie było
ponownie zadawane temu samemu pracownikowi po medianie \GapMedian{} dni
(IQR \GapQOne--\GapQThree).

\subsection{Dane}
Dane zostały zorganizoane w firmie listy zawierającej  moduł, pytanie, wybraną opcję,
kod odpowiadającego, znacznik czasu oraz poprawność odpowiedzi (flaga true/false). Wykluczono  \NExcludedItems{} pytania, w których zarejestrowano
tylko jedną opcję odpowiedzi (\NExcludedItemResp{} odpowiedzi), ponieważ nie niosą one informacji
o błędach. Zbiór analityczny obejmował \NResponses{} odpowiedzi (tab.~\ref{tab:zbior_danych}).
Dane osób odpowowiadających zostały przekazane w formie metadanych, zawierające informacje o płci panelisty oraz przedziale wiekowym (co 10 lat).

\fillin{Podstawa prawna przetwarzania danych spseudonimizowanych (wiek, płeć) oraz zgoda komisji
bioetycznej lub zwolnienie z obowiązku jej uzyskania, jeśli wymaga tego czasopismo.}

\subsection{Definicje}
\emph{Sekwencja} to uporządkowana chronologicznie seria prób jednego pracownika dotyczących jednego
pytania; \emph{numer próby} to pozycja odpowiedzi w jej sekwencji (przy jednakowym czasie o kolejności
decydował identyfikator rekordu). \emph{Poprawność} to odsetek poprawnych odpowiedzi. \emph{Odstęp}
to czas w dniach od poprzedniej próby w tej samej sekwencji. \emph{Staż} to czas od pierwszej
odpowiedzi pracownika w programie. Zmianę na poziomie pracownika zdefiniowano jako różnicę odsetka
poprawnych odpowiedzi między próbą~3 a próbą~1, obliczoną na podstawie własnych sekwencji pracownika
liczących co najmniej trzy próby.

\subsection{Analiza statystyczna}
Odsetki poprawnych odpowiedzi w statystykach opisowych podano z 95-procentowymi przedziałami ufności (PU) z bootstrapu
klastrowego, w którym losowano pracowników ze zwracaniem (\BootReps{} replikacji); procedura ta
uwzględnia zależność odpowiedzi udzielanych przez tę samą osobę. Ponieważ osoby kontynuujące do kolejnych prób mogą różnić się od
tych, które przestają, zagregowane krzywe uczenia mogą być zniekształcone przez ubywanie obserwacji
\citep{murray2013,goutte2018}; dlatego krzywe uczenia obliczono dodatkowo dla zbalansowanej kohorty
sekwencji liczących co najmniej \CohortK{} prób. Zależność odsetka błędów od numeru próby opisano
funkcją potęgową \citep{newell1981}. Ponieważ harmonogram powtórek nie był adaptacyjny, sprawdzono,
czy ponowne zadanie pytania zależało od wyniku poprzedniej próby.

Poprawność odpowiedzi modelowano regresją logistyczną z efektami stałymi pracownika i pytania oraz
błędami standardowymi klastrowanymi względem pracowników, tak aby efekt numeru próby był szacowany
w obrębie osób i w obrębie pytań. Model~1 zawierał numer próby w skali log$_2$, dzięki czemu iloraz
szans (OR) odnosi się do podwojenia liczby prób. Model~2, dopasowany do prób powtórnych, uwzględniał
dodatkowo wynik poprzedniej próby, odstęp od poprzedniej próby (log$_2$ dni) oraz staż, analogicznie
do modeli \emph{performance factors analysis} i \emph{half-life regression} stosowanych do logów
uczących się \citep{pavlik2009,settles2016}. Model~3, dopasowany wyłącznie do pierwszych prób,
weryfikował, czy poprawność odpowiedzi na nowo napotkane pytania zmieniała się wraz ze stażem (PB4).
W analizach wrażliwości wykorzystano zbalansowaną kohortę, obcięcie do 10 prób, warstwy według długości
sekwencji (kontrola obciążenia typu \emph{pattern-mixture}), wykluczenie pracowników o niskiej
aktywności, korektę o staż, uogólnione równania estymujące (GEE) z wymienną macierzą korelacji
roboczej \citep{liang1986} oraz model mieszany ze skrzyżowanymi losowymi wyrazami wolnymi pracowników
i pytań \citep{baayen2008}, na podstawie którego wyznaczono współczynniki korelacji wewnątrzklasowej
(ICC) w skali latentnej oraz brzegowy i warunkowy współczynnik $R^2$ \citep{nakagawa2013}.

Przejścia między kolejnymi próbami zestawiono w tabeli kontyngencji. Dla dwóch następujących po sobie błędów porównano odsetek ponownych wyborów
tego samego dystraktora z odsetkiem oczekiwanym przy losowym wyborze spośród zaobserwowanych
dystraktorów oraz przy wyborze proporcjonalnym do popularności każdego dystraktora w danym pytaniu.
Zmianę na poziomie pracowników testowano testem znaków i testem rangowanych znaków Wilcoxona
(korelacja rangowo-dwuseryjna jako miara wielkości efektu). Pytania opisano wskaźnikami klasycznej
teorii testu obliczonymi na pierwszych próbach: trudnością (odsetek poprawnych odpowiedzi),
dyskryminacją (korelacja pytanie--reszta, tj.\ z poprawnością pracownika w pozostałych pytaniach) oraz
funkcjonowaniem dystraktorów; dystraktory wybierane przez mniej niż 5\% odpowiadających uznano za
niefunkcjonalne \citep{tarrant2009,gierl2017}. Trendy poprawności w kolejnych próbach dla
poszczególnych pytań testowano regresją logistyczną z błędami standardowymi klastrowanymi względem
pracowników, kontrolując odsetek fałszywych odkryć (FDR) na poziomie \Alpha{} \citep{benjamini1995}.
Analizy wykonano w języku Python \VerPython{} z użyciem bibliotek pandas \VerPandas, NumPy \VerNumpy,
SciPy \VerScipy, statsmodels \VerStatsmodels{} i Matplotlib \VerMatplotlib{}.
%  (ziarno generatoraliczb losowych \Seed).
